\documentclass{article}
\usepackage[utf8]{inputenc}
\usepackage{amsmath, amssymb}

\begin{document}
	
	\section{Cuestión 2}
	\begin{center}
		$G_v = \frac{V_{R_L}}{V_g}$; $R_1 = 1\text{ k}\Omega$, sobre una resistencia de carga $R_L = 10\text{ k}\Omega$.
	\end{center}
	
	
	\[
	\left.
	\begin{aligned}
		v^+ &= v^- \\
		v^+ &= 0 \\
		\frac{V_e - v^-}{R_1} &= \frac{v^- - v_s}{R_2}
	\end{aligned}
	\right\}
	\implies \frac{V_e}{R_1} = -\frac{v_s}{R_2} \implies G_v = -\frac{R_2}{R_1}
	\]
	
	\begin{center}
		\text{Luego}
	\end{center}		
	\[ 	
	 \quad R_2 = -G_v \cdot R_1\implies R_2 = 10 \cdot 1000 \implies R_2 = 10k\Omega
	\]
	
\end{document}